\documentclass[12pt]{article}
\usepackage[utf8]{inputenc}
\usepackage{amsmath}
\usepackage{graphicx}
\usepackage{natbib}
\usepackage{hyperref}
\usepackage{lipsum} % For dummy text, you can remove this in your final document

\title{Bioinspired Design of Materials for Energy Storage: Principles and Applications in Batteries and Supercapacitors}
\author{Your Name}
\date{\today}

\begin{document}

\maketitle

\begin{abstract}
Bioinspired design offers innovative strategies for enhancing energy storage materials by mimicking natural structures and processes. This paper explores bioinspired design principles, evaluates material properties in batteries and supercapacitors, assesses performance metrics, and discusses future research directions in the field.
\end{abstract}

\section{Introduction}
Bioinspired design involves emulating structures and processes from nature to create advanced materials for energy storage applications. This paper reviews the principles of bioinspired design in materials science, discusses its relevance to improving energy storage technologies, and outlines the structure of batteries and supercapacitors as key focus areas.

\section{Design Principles}
\subsection{Structural Mimicry}
Natural structures, such as hierarchical arrangements found in bones and trees, inspire designs for electrodes and electrolytes in batteries and supercapacitors \citep{lee2012bioinspired, liu2019bioinspired}.

\subsection{Functional Mimicry}
Biochemical processes, like ion transport in biological membranes, guide the development of ion-conductive materials for enhancing energy storage device performance \citep{luo2016bioinspired, jiang2020bioinspired}.

\section{Material Properties}
\subsection{Electrode Materials}
Bioinspired electrode materials exhibit enhanced surface area, conductivity, and ion diffusion properties crucial for efficient charge storage and cycling stability \citep{li2017bioinspired, wang2018bioinspired}.

\subsection{Electrolytes}
Natural electrolyte analogs improve ion transport kinetics and electrochemical stability in batteries and supercapacitors, reducing internal resistance and enhancing device longevity \citep{jin2015bioinspired, guo2017bioinspired}.

\section{Performance Evaluation}
\subsection{Energy Density and Power Density}
Bioinspired materials enhance energy and power densities by optimizing electrode-electrolyte interfaces and charge transfer mechanisms, improving overall device efficiency \citep{liu2020bioinspired, zhang2021bioinspired}.

\subsection{Cycling Stability}
Improved material durability and resistance to degradation contribute to extended cycling lifetimes and reliability of energy storage devices \citep{cao2018bioinspired, yang2019bioinspired}.

\section{Future Research}
Future directions include advancing bioinspired designs for multi-functional materials, integrating renewable resources, and optimizing scalability and manufacturing processes to realize practical applications in energy storage technologies.

\section*{References}
\bibliographystyle{plainnat}
\bibliography{prompt7_35}

\end{document}
